\section{Threat Model}
\label{sec:threat}
% =====================================================================

\paragraph{System.}
An LLM agent $A$ controls a CPS plant $P$ (water-treatment substrate
SWaT P1+P2, with raw-water intake T101, chemical-dosing buffer T201,
transfer valves $\{$MV101, MV201$\}$, raw-water pumps
$\{$P101, P102$\}$, dosing pumps $\{$P201--P206$\}$, and analyser
sensors $\{$LIT101, LIT201, FIT101, FIT201, AIT201--AIT203$\}$). The
agent issues actions to $P$ via Model Context
Protocol~\cite{mcp_spec} \emph{tool calls}: each tool exposes a
natural-language description, an input schema, and an
implementation that wraps a PLC-style command on $P$. The defender
places a physics-grounded verifier $V$ between the MCP server and
$A$'s tool catalog.

\paragraph{Attacker.}
The attacker controls the natural-language description of one or more
MCP tools, the implementation behind those tools, or both, and can
optionally pre-set the physical actuator state on $P$ before the
agent's session begins. The attacker may additionally spoof sensor
reads at the network layer (post-admission). The attacker cannot
modify the verifier $V$, the digital twin $\mathit{DT}$, the lifter
$L_{\text{verify}}$, or the formal matcher; nor can the attacker
bypass the admission gate.

\paragraph{Attack classes.}
We organise attacks into three categories aligned with the empirical
campaigns.
\begin{enumerate}
\item \textbf{Description-poisoning} (A1, A2, A3): the attacker writes
      a description that mismatches the tool's implementation. A1 is
      type-confusion (the description claims a tool drains while the
      implementation fills); A2 is magnitude poisoning on a parametric
      tool (the description names a unit/scale that the
      implementation does not honour); A3 is sensor aliasing (the
      description claims one sensor while the implementation returns
      another).
\item \textbf{Composition} (Z1, Z2, Z3): the attack manifests through
      the joint behaviour of multiple individually-admitted tools.
      Z1 is a single-tool transient overshoot exploiting steady-state
      semantics in the lifted band; Z2 is overdose magnitude poisoning
      (mirror of A2); Z3 is the canonical strict-composition witness
      (described in detail in \S\ref{sec:t3-witnesses}).
\item \textbf{Runtime} (W1, W2): the attack is invisible at admission
      time and manifests after the agent's session begins. W1 is a
      sub-noise-floor stealth drain that respects the description's
      claimed state-zero invariant within per-sample fidelity but
      accumulates a cumulative drift exceeding the matcher's
      cumulative budget. W2 is post-admission sensor spoofing of
      LIT101 while pumps drive the tank to overflow.
\end{enumerate}

\paragraph{Trusted computing base.}
Table~\ref{tbl:tcb} summarises the trusted computing base. PG-DSL is
only as good as DT fidelity, grammar coverage, and the lifter's
bounded semantic error on $\mathcal{C}$.

\begin{table}[t]
\centering
\footnotesize
\caption{Trusted computing base for PG-DSL. PG-DSL provides no
defence if a trusted component is itself compromised.}
\label{tbl:tcb}
\begin{tabular}{@{}p{0.30\columnwidth}p{0.18\columnwidth}p{0.40\columnwidth}@{}}
\toprule
Component & Trusted? & Rationale \\
\midrule
MCP description & No & attacker-controlled \\
Tool impl. & Partial & may be poisoned; DT observes \\
Digital twin & Yes & verifier substrate, fidelity-bounded \\
Lifter & Yes, fallible & bounded $\varepsilon_L$; T2 quantifies \\
Formal matcher & Yes & deterministic check \\
Runtime telemetry & Partial & W2 spoofs post-admission reads \\
Runtime IDS & Separate & T3 composition partner \\
\bottomrule
\end{tabular}
\end{table}

\paragraph{Out of scope.}
Network-level attacks against the verifier itself, hardware ageing,
chemistry kinetics beyond linearised conductivity, and side-channel
attacks on $L_{\text{verify}}$ (e.g., adversarial probes designed
specifically against the lifter's training distribution) are out of
scope for this paper. These are deferred to T3's runtime composition
partner or to future work.

\paragraph{Real-agent threat realisation.}
A bounded real-agent campaign on the SWaT P1+P2 substrate validates
the threat-model claim end-to-end: a \texttt{qwen3:14b} agent
operating against a real-MCP server via the official Python MCP SDK
falls to description-channel poisoning (A1, type-confusion overflow)
in $10/10$ replications without defence; remains vulnerable to
post-admission sensor spoofing (W2) in $9/10$ replications regardless
of admission-time defence (W2 is a canonical $A_{\text{runtime}}$
witness whose visibility surface is runtime telemetry rather than
admission-time descriptions). With
canonical PG-DSL admission applied, A1 ASR drops to the with-defence
row of Table~\ref{tbl:real-agent-asr} --- the admission gate excludes
the poisoned tool from the agent's catalog before the agent is ever
exposed to the description --- while W2 ASR is unchanged, as
predicted by the T3 partition (W2 is a canonical $A_{\text{runtime}}$
witness). The campaign uses $N=10$ reps $\times$ $30$-step horizons
$\times$ four conditions $\{A_1, A_2, W_2, A_2^{\text{eng}}\}$ at
$10$~s simulated plant-tick per agent action; reps sweep across
realistic operating-condition variation (see \S\ref{sec:eval-real-agent}),
not stochastic LLM-policy independence. Full per-rep traces and
aggregate Wilson~95\% CIs are in the artifact. Detailed numbers
appear in \S\ref{sec:eval-real-agent}.

% =====================================================================
